\section{One-body magic cannot bound the determinant support}
\label{app:theorem}
\begin{proposition}[Coefficient-insensitivity obstruction]
\label{prop:obstruction}
Let $\mu$ be any functional of a fixed finite set of reduced density matrices, continuous in the
state. Then $\mu$ cannot usefully upper-bound the \emph{raw} determinant support (the unthresholded
Slater rank): for every $\epsilon>0$ there are states with $|\mu-\mu(\ket{D_0})|<\epsilon$ (for
$\mathcal F_k$, $\mu<\epsilon$) whose raw support is every determinant of the particle-number sector, so
any bound $|\mathcal S|_{\rm raw}\le f(\mu)$ valid on the sector takes the trivial value, the sector
dimension, at values of $\mu$ arbitrarily close to $\mu(\ket{D_0})$. The count is lower- but not
upper-semicontinuous; the CP tensor rank fails upper semicontinuity in the same way and, unlike
the count, is not even lower-semicontinuous: its sets of rank at most $r$ are not
closed~\cite{desilvalim2008} (that reference names the two semicontinuities the other way round,
calling matrix rank upper-semicontinuous). Nor can the invariant $\mathcal F_1$ usefully lower-bound it:
for $N$ electrons on $2n_{\rm o}$ modes it implies $|\mathcal S|_{\rm raw}\ge N/(N-\mathcal F_1/4)$,
which never exceeds $2n_{\rm o}/N$, and at half filling ($N=n_{\rm o}$), where that is $2$, any
increasing $f$ with $|\mathcal S|_{\rm raw}\ge f(\mathcal F_1)$ satisfies $f\le2$.
\end{proposition}
\begin{proof}
Rank is a coefficient-insensitive count. Take
$\ket{\Psi_\epsilon}=\ket{D_0}+\epsilon\sum_{k=1}^{K}\ket{D_k}$ with the $D_k$ distinct determinants of
the sector of $\ket{D_0}$. As $\epsilon\to0$ every reduced density matrix converges to that of
$\ket{D_0}$, so any continuous $\mu\to\mu(\ket{D_0})$ (for $\mathcal F_k$, to $0$), while the raw
support stays $K+1$ for every $K$ up to the sector dimension minus one. For the second part, every
natural orbital lies in the span of the at most $N|\mathcal S|_{\rm raw}$ orbitals occupied in the
support, so $\mathrm{rank}\,\gamma\le N|\mathcal S|_{\rm raw}$; with $\mathrm{tr}\,\gamma=N$ and
$\mathrm{tr}\,\gamma^{2}=N-\mathcal F_1/4$, the Cauchy--Schwarz inequality
$(\mathrm{tr}\,\gamma)^{2}\le\mathrm{rank}\,\gamma\,\mathrm{tr}\,\gamma^{2}$ gives the bound, and
$\mathrm{rank}\,\gamma\le2n_{\rm o}$ caps it at $2n_{\rm o}/N$. At half filling with $n_{\rm o}\ge2$ the two-determinant cat
$(\ket{1\cdots10\cdots0}+\ket{0\cdots01\cdots1})/\sqrt2$, whose determinants differ in $N\ge2$ orbitals,
has $\gamma=\tfrac12\openone$, hence $\mathcal F_1=2n_{\rm o}$ (maximal) and Slater rank $2$, so
$f(2n_{\rm o})\le2$; for $n_{\rm o}=1$ every state is a single particle on two modes, of raw support at
most $2$.
\end{proof}
\noindent The obstruction is about the unthresholded rank; the operative, coefficient-\emph{sensitive}
support $|\mathcal S|_\varepsilon$ of Sec.~\ref{sec:nogo} collapses to $1$ along the same
$\epsilon\to0$ path, so continuity alone does not obstruct it. What obstructs $|\mathcal S|_\varepsilon$
itself is the orbital-rotation symmetry---no invariant functional survives it---together with the two
equal-weight witnesses below, whose amplitudes are $\Theta(1)$ and therefore threshold-robust.
Two explicit witnesses rule out the natural repairs. The two-determinant cat of
Prop.~\ref{prop:obstruction} is maximal in $\mathcal F_1$ yet of Slater rank $2$, so at half filling no
$\mathcal F_1$-monotone bounds $|\mathcal S|$ from below beyond $2$; a paired (BCS) Gaussian state has
Gaussian rank $1$ yet fractional occupations, so no occupation functional bounds it from above. The
computable lower bounds on $|\mathcal S|$ that are not capped in this way are the basis-dependent
orbital-Schmidt (matrix-flattening) ranks~\cite{schollwock2011}---the tensor-network bond
dimension---to which any Gaussian-invariant magic measure is blind. That lower bound is exact:

\begin{proposition}[The bond dimension lower-bounds the support]
\label{prop:chibound}
For any state and any bipartition of the spin-orbitals into $(\mathrm L,\mathrm R)$, the Schmidt rank
$\chi$ across the cut (the minimal matrix-product-state bond dimension there) satisfies
\begin{equation}
\chi \;\le\; |\mathcal S|,
\end{equation}
where $|\mathcal S|$ is the number of fixed-basis Slater determinants of nonzero amplitude. Hence the
minimal bond dimension is a rigorous, basis-dependent lower bound on the determinant support. The bound is
one-sided: the geminal witness of Thm.~\ref{thm:witness} has $\chi=2$ yet $|\mathcal S|_{\mathrm{pair}}=2^{K}$,
so $|\mathcal S|$ can exceed $\chi$ exponentially already in the pairing basis, which for that family is
a natural-orbital basis, while the orbital-invariant one-body magic $\mathcal F_1$ cannot usefully
upper-bound $|\mathcal S|$ and at half filling cannot lower-bound it beyond $2$
(Prop.~\ref{prop:obstruction}).
\end{proposition}
\begin{proof}
Write $\ket{\Psi}=\sum_{D\in\mathcal S}c_D\ket{D}$ over its support. Each determinant factorizes across the
cut, $\ket{D}=\ket{D_{\mathrm L}}\otimes\ket{D_{\mathrm R}}$, so the matricization
$A_{D_{\mathrm L},D_{\mathrm R}}=c_D$ has at most $|\mathcal S|$ nonzero entries; its rank---the Schmidt
rank $\chi$---is therefore at most $|\mathcal S|$.
\end{proof}
\noindent Since a matrix-product-state simulation costs $\mathrm{poly}(\chi)$, Prop.~\ref{prop:chibound}
says that \emph{in the sampler's own basis a tensor network is never less compact than the sampled
subspace}: this is why we make no advantage claim in one dimension, where $\chi$ is small, and why
any advantage must live in the higher-entanglement regime where $\chi$ itself---and with it the classical
tensor-network cost---grows (the chain-to-ladder jump of Fig.~\ref{fig:ladder}). Empirically $\chi$ is also the
operative correlate of the \emph{measured} cost: across the $74$ exact ground states of
Fig.~\ref{fig:master} the thresholded support $|\mathcal S|_\varepsilon$ tracks $\chi$ (pooled Spearman
$\rho=0.72$; $\rho=0.90$ within the $n=38$ molecular suite, where however one-body magic scores a
comparable $\rho=0.86$---the two resources are separated not there but by the Hubbard sweep, where
$\mathcal F_1$ anti-orders the support in all six strata while $\chi$ orders it at a mean $\rho$ of
$0.948$; the pooled coefficients for that sweep, $\rho=0.57$ and $-0.11$, are Simpson reversals of the
six strata and are not the evidence here, Sec.~\ref{sec:nogo:chi}). Two caveats belong with that number. A rank correlation is not a statement of tightness,
and per the geminal witness the bound can be exponentially loose. And the proposition is proved for the
raw support: once amplitudes are thresholded the inequality need not survive---here $|\mathcal S|$ is cut
at $99\%$ of the weight while $\chi$ is cut at $\varepsilon=0.05$, a tighter tolerance---and
$|\mathcal S|_\varepsilon$ falls below $\chi$ in $7$ of the $38$ molecular points. What stands is that $\chi$ is a rigorous floor on
the raw support and an empirical correlate of the measured one, whereas $\mathcal F_1$ orders the
support with a sign that is not stable from family to family (Sec.~\ref{sec:nogo:chi}).

